%===============================================================================
\section{Experimental Setup}\label{sec:setup}
%===============================================================================

\subsection{Data sets and rule-based field roles}\label{sec:data}
Table~\ref{tab:data} summarizes the four data sets. In all of them, candidate fields are
information that the disclosure rules~\cite{nea2024disclosure} release \emph{before} the
target, and targets are either specific information (unit output, Art.~17(3)), quantities not on
the disclosure list at hourly resolution (hourly line flow; only daily averages of important
lines are listed), or actual operating quantities released on the following day (actual load,
total generation, tie-line exchange; Appendix items 6.57--6.66). Using the release timing to
assign roles does not involve any temporal split of the data; all splits are random.

\emph{RTS-GMLC.} The IEEE Reliability Test System 2019 update~\cite{barrows2020rts} provides a
73-bus network with one year of realistic hourly load, wind, solar and hydro series. We clear
an hourly DC optimal power flow with piecewise-linear offers built from the published heat-rate
curves (8,784 hours, 26.7\% congested). Candidates are the 23 fields a market would publish
early (zonal load forecasts, renewable and hydro forecasts, reserve requirements, system energy
price and nine nodal prices, three congestion indicators); targets are the outputs of a coal and
a combined-cycle unit and the hourly flow on internal line C35.

\emph{NEM.} One year (2024) of the Australian National Electricity Market~\cite{aemo2024nemweb},
averaged to hours. Candidates are regional demand, available capacity, semi-scheduled wind and
solar forecasts, regional prices and interconnector flows (30 fields); targets are the actual
outputs of a coal unit (BW01), a hydro unit (TUMUT3) and a gas combined-cycle unit (PPCCGT).
Unit output is public in the NEM but is specific information under the Chinese rules; this data
set therefore contains a genuinely sensitive field observed in the real world.

\emph{PJM and CAISO.} System-level hourly data for one year. Targets are the actual quantities
released on the next day: PJM metered load, total generation and net actual interchange, and
CAISO actual load. Candidates are forecasts, day-ahead and real-time prices and their
components, ancillary-service requirements and cleared volumes, and schedules (22--35 fields).
Fields released on the same day as the target (fuel-mix generation, actual renewable output,
other balancing areas' actual load) are excluded, as are the load forecasts of the target area,
which are ex-ante copies of the target (Section~\ref{sec:problem}). Near-copies that are distinct
quantities (e.g.\ sub-area forecasts for CAISO load) are kept; removing them is analysed as a
sensitivity case.

\begin{table}[!tbp]
\centering\small
\caption{Data sets. $n$: train/test rows; $p$: candidate fields; sets: field sets with exact
three-attacker ground truth.}\label{tab:data}
\resizebox{\textwidth}{!}{%
\begin{tabular}{llcccl}
\toprule
Data set & Targets (rule basis) & $n$ & $p$ & sets & max size \\
\midrule
RTS-GMLC & unit 223, unit 221 (specific info.); line C35 (not listed) & 6149/2635 & 23 & 10,902 & 4\\
NEM 2024 & units BW01, TUMUT3, PPCCGT (specific info.) & 6149/2635 & 30 & 31,930 & 4\\
PJM (load) & metered load (next-day) & 4589/1967 & 22 & 9,108 & 4\\
PJM (gen./interch.) & total generation, net interchange (next-day) & 4589/1967 & 24 & 12,950 & 4\\
CAISO & actual load (next-day) & 6105/2617 & 35 & 59,535 & 4\\
\bottomrule
\end{tabular}}
\end{table}

\subsection{Ground truth}
All variables are standardized with training statistics; 70/30 random train/test splits are
used, with 15\% of the training rows as validation. The attacker family has three members: a
single-target multilayer perceptron (two hidden layers of 128 units, dropout 0.15, Adam,
early stopping on validation), a multi-target perceptron trained on all targets of a data set,
and histogram gradient boosting with two configurations. For every set of up to four fields,
every member is trained on that set alone; the member and epoch are chosen on validation and the
test $\Rtwo$ is reported, followed by the monotone closure~\eqref{eq:closure}. This amounts to
124,425 field sets and several hundred thousand trained models. Neither the
estimator nor any baseline has access to these values; they are used only for evaluation.

\subsection{Baselines}
Field scores compared with $M^{(2)}$ are: Pearson $r^2$, squared Spearman correlation, mutual
information~\cite{kraskov2004mi}, distance correlation~\cite{szekely2007dcor}, single-field
retraining $M^{(0)}$ (current practice), LOCO~\cite{lei2018loco}, permutation
importance~\cite{breiman2001rf}, SAGE with marginal imputation~\cite{covert2020sage}, and an
inference-graph score that follows the propagation idea of~\cite{hu2026ignn} without its expert
inputs: edge weights are the single-field predictability of one field from another (gradient
boosting, validation $\Rtwo$), the target node carries risk one, and risk is propagated over three
hops with decay 0.5. All baseline scores are computed on training rows only.

Estimators compared with ours are: linear and quadratic readouts without a backbone, an in-set
quadratic dictionary (values, squares and pairwise products of the fields in $S$), the shared
output head of the same backbone, an untrained (randomly initialized) backbone with the same
readout, and warm-start fine-tuning of the whole network for each set~\cite{sun2025warmstart}.

\subsection{Metrics}
For set values: mean absolute error of $\hat V$. For field risk: mean absolute error and bias of
$\hat M^{(2)}$, Spearman correlation across fields, and the certified ratio
$\sum_i L^{(2)}_{i\to y}/\sum_i M^{(2)}_{i\to y}$ when $k$ backgrounds per field are retrained.
For decisions at thresholds $\tau\in\{0.5,0.7\}$: recall and precision of critical fields; the
number of fields withheld and the number of minimal unsafe sets of size $\le 3$ left unbroken.
